\section{Prerequisiti}

Algebra: calcolo combinatorio, teoria degli insiemi, relazioni e funzioni, chiusura dell'anello dei polinomi rispetto alla composizione (che rende particolarmente interessanti gli algoritmi polinomiali).\\

Elementi di informatica teorica: macchine di Turing, linguaggi ricorsivi e ricorsivamente enumerabili.\\
Puó aiutare ricordare il funzionamento degli automi a stati finiti non deterministici per comprendere le Macchine di Turing non deterministiche.\\
Puó aiutare ricordare il problema della fermata per la dimostrazione dell'Hierarchy Theorem (\ref{sec:hierarchy_theorem}).\\

Algoritmi e strutture dati: notazione asintotica, dimostrazioni di correttezza, analisi di complessitá degli algoritmi, teoria dei grafi.\\

Logica: sintassi e semantica della logica proposizionale e del primo ordine, in particolare la definizione di soddisfacibilitá, fondamentale per $\SAT$ (\ref{problema:sat}). \\
Puó essere utile confrontare la descrizione di un circuito hamiltoniano in logica proposizionale con la riduzione di $\HAMILTONPATH$ a $\SAT$ (\ref{riduzione:hamiltonpath_sat}).\\
Puó aiutare ricordare le regole del calcolo deduttivo e alcune tautologie per seguire meglio alcune dimostrazioni.\\